\ExerciseSection

\begin{enumerate}
\def\labelenumi{\arabic{enumi})}
\item
  设 \(a\) 是一个非零复数，\(n\) 是一个 \(>0\) 的整数。对每个实数 \(r > 0\) （满足 \(r^n \leqslant |a|\) ），证明存在 \(z \in \mathbf{C}\) 使得 \(|z| = r\) 且 \(|a + z^n| = |a| - r^n\) 。由此推出：若 \(f\) 是次数 \(>0\) 的复系数多项式，则不等式 \(|f(z_0)| \leqslant |f(z)|\) 不可能对一个邻域中的所有点 \(z\) 成立，该邻域属于某点 \(z_0\) ，只要 \(f(z_0) \neq 0\) 。
\item
  证明：若 \(f\) 是任一复系数多项式且不恒等于零，则存在实数 \(r > 0\) 使得
\end{enumerate}

\[
\left| f (z) \right| > \left| f (\mathbf {o}) \right|
\]

对一切满足 \(|z| \geqslant r\) 的 z 成立。于是，借助习题 1 与魏尔斯特拉斯定理（第四章，§ 6，第 1 号，定理 1），给出 \(\mathbf{C}\) 是代数闭域这一事实的另一个证明（考察函数 \(f\) 在由点 \(z\) （满足 \(|z| \leqslant r\) ）组成的紧集上的情况）。

\begin{enumerate}
\def\labelenumi{\arabic{enumi})}
\setcounter{enumi}{2}
\item
  不利用定理 1 证明：映射 \(z \to z / \overline{z}\) 是拓扑群 \(\mathbf{C}^*\) 到拓扑群 \(\mathbf{U}\) 上的一个严格态射。由此推出：映射 \(t \to (\mathrm{i} + it) / (\mathrm{i} - it) \left[ (\mathrm{i} + it) / (\mathrm{i} - it) = -\mathrm{i} \right.\) （若 \(t = \infty\) ）是拓扑群 \(\tilde{\mathbf{R}}\) （第 6 号）到拓扑群 \(\mathbf{U}\) 上的一个同构，从而给出定理 1 的另一个证明。
\item
  设 K 是 R 的最小毕达哥拉斯子域，并设 \(\mathrm{K}^{\prime}=\mathrm{K}(i)\) 。设 G 是 \(K^{\prime}\) 中绝对值为 1 的元素所成的乘法群（U 的一个子群）。证明 G 不同构于 K 中数 mod 1 的加法群。（注意在后一群中存在任意素数阶 p 的元素；另一方面，若 p 是使 p-1 不是 2 的方幂的素数，则注意 \(K^{\prime}\) 的每个元素在 Q 上的次数都是 2 的方幂，由此证明 G 中除 1 外不含 p 次单位根。）
\end{enumerate}

